\documentclass[11pt]{article}
\usepackage[margin=0.78in]{geometry}
\usepackage{amsmath,amssymb,amsthm,booktabs,longtable,array}
\setlength{\parindent}{0pt}
\setlength{\parskip}{0.45em}
\newcommand{\Q}{\mathbb Q}
\newcommand{\om}{\omega _0}
\newcommand{\rc}{\rho _c}
\newcommand{\rs}{\rho _*}
\newcommand{\sgn}{\operatorname{sgn}}
\newtheorem{lemma}{Lemma}[section]

\title{Hand-visible analytic derivations for 103 finite-ledger predicates}
\author{}
\date{}

\begin{document}
\maketitle

\section{Scope and logical convention}

This note exposes, in hand-checkable form, the predicates not assigned to
the other five ledger packets.  They are precisely
\[
  11\ \hbox{threshold predicates}
  +48\ \hbox{zero predicates}
  +43\ \hbox{lower-closure predicates}
  +1\ \hbox{cap--74 predicate}=103.
\]
The omitted lower-closure predicate is the single six-cap tuple identity;
it belongs to the cap-inventory packet.  Every inequality below is proved
from a displayed analytic inequality, integer arithmetic, rational cross
multiplication, or positive-side squaring.  An executable ledger may check
the same arithmetic, but no executable output is used as a premise.

We use
\[
 \om=\frac{\sqrt3}{2\pi}-\frac16,
 \quad C_0=1+\frac{8\sqrt2}{15},
 \quad \rs=\frac{\om}{2C_0},
 \quad \rho _0=\frac1{\sqrt{337}},
 \quad \sigma=\frac{3\om}{4},
 \quad \rc=\frac1{1+2\pi},
 \quad d=\frac{\sqrt{337}}2.
\]

The elementary enclosure
\begin{equation}
 \frac{333}{106}<\pi<\frac{355}{113}<\frac{22}{7}
 \label{eq:pi-bounds}
\end{equation}
is analytic: the lower and middle bounds follow from alternating truncations
of Machin's identity
\(\pi=16\arctan(1/5)-4\arctan(1/239)\), and the last one is immediate by
cross multiplication.  For reference, \(355/113\) exceeds the upper
Machin truncation by the positive rational numerator
\(45167474711\) over denominator
\(189820334689453125\).  The classical integral
\[
 0<\int_0^1\frac{x^4(1-x)^4}{1+x^2}\,dx=\frac{22}{7}-\pi
\]
also gives the final upper bound directly.

\section{The eleven threshold predicates}

First we obtain a rational two-sided enclosure for \(\om\).  The exact
integer residuals
\[
 3\,1101^2 113^2-607^2 355^2=1998482>0,
 \qquad
 25\,333^2-243\,106^2=41877>0
\]
and \eqref{eq:pi-bounds} imply
\[
 \sqrt3>\frac{607\,355}{1101\,113}>\frac{607\pi}{1101},
 \qquad 9\sqrt3<5\pi.
\]
Consequently
\begin{equation}
 \frac{40}{367}<\om<\frac19.
 \label{eq:omega-bounds}
\end{equation}
Indeed, \(607/2202-1/6=40/367\), while
\(5/18-1/6=1/9\).

Since \(2\,32^2-45^2=23>0\), one has
\(\sqrt2>45/32\), hence \(C_0>7/4\) and
\(\rs<2/63\).  The following chain consists entirely of strict rational
comparisons, with the indicated radical comparisons justified by squaring:
\[
 \rs<\frac2{63}<\frac1{\sqrt{337}}<\frac1{18}
 <\frac3{40}<\frac{3\om}{4}<\frac1{12}<\frac7{51}
 <\frac1{1+2\pi}<\frac17<\frac{39}{50}<\frac78.
\]
Here
\[
 4\cdot337<63^2,\qquad 18^2<337,\qquad 367<400,
 \qquad 51<84,\qquad 367<420,
\]
and \(7/51<\rc<1/7\) follows respectively from
\(\pi<22/7\) and \(\pi>3\).  This proves the five ratio-order rows.
It also proves
\[
 \om<\frac19<\frac7{51}<\rc.
\]
Moreover
\[
 \om d>\frac{20\sqrt{337}}{367}>1,
 \qquad 400\cdot337-367^2=111>0.
\]

The four terminal comparisons follow from the same bounds:
\begin{align*}
 \sqrt{114}&<18<\frac2\om,\qquad
 &\frac2\om&<\frac{367}{20},\\
 \rc&<\frac17<\frac{60}{367}<\frac{3\om}{2},
 &\om\sqrt{114}&>\frac{40\sqrt{114}}{367}>1.
\end{align*}
The last positive-side squaring has residual
\[
 1600\cdot114-367^2=47711>0.
\]

Thus the one-to-one threshold registry is
\begin{center}
\small
\begin{tabular}{ccl}
\toprule
ID&source row&proved predicate\\
\midrule
T1--T5&1--5&\(\rs<\rho _0<\sigma<\rc<39/50<7/8\), one adjacent pair per row\\
T6&6&\(\om<\rc\)\\
T7&7&\(\om d>1\)\\
T8&31&\(\sqrt{114}<2/\om\)\\
T9&32&\(2/\om<367/20\)\\
T10&33&\(\rc<3\om/2\)\\
T11&34&\(\om\sqrt{114}>1\)\\
\bottomrule
\end{tabular}
\end{center}

\section{The forty-eight zero specializations}

\subsection{A recurrence and elementary tangent bounds}

Put
\[
 \mathsf j_\ell(x)=\sqrt{\frac{\pi}{2x}}J_{\ell+1/2}(x),
 \qquad P_\ell(x)=x^{\ell+1}\mathsf j_\ell(x).
\]
Then
\begin{equation}
 P_{\ell+1}=(2\ell+1)P_\ell-x^2P_{\ell-1},
 \qquad P_0=\sin x,\qquad P_1=\sin x-x\cos x.
 \label{eq:P-rec}
\end{equation}
Write
\(P_\ell=A_\ell(t)\sin x+xC_\ell(t)\cos x\), \(t=x^2\).
The recurrence gives
\begin{align*}
 A_2&=3-t,&C_2&=-3,\\
 A_3&=15-6t,&C_3&=t-15,\\
 A_5&=945-420t+15t^2,&C_5&=-945+105t-t^2,\\
 A_6&=10395-4725t+210t^2-t^3,
 &C_6&=-10395+1260t-21t^2,\\
 A_7&=135135-62370t+3150t^2-28t^3,
 &C_7&=-135135+17325t-378t^2+t^3.
\end{align*}
For \(0<s<\pi/2\), differentiation and integration give
\begin{equation}
 \tan s>s,\qquad \tan s>s+\frac{s^3}{3}.
 \label{eq:tan-lower}
\end{equation}
For \(0<s<1\), \(\sin s<s\) and
\(\cos s>1-s^2/2\) give
\begin{equation}
 \tan s<\frac{s}{1-s^2/2}.
 \label{eq:tan-upper}
\end{equation}

\subsection{Sixteen half-period locations}

For each internal face, the first displayed margin below is
\(w-(h_L/2)(22/7)\); the second is
\((h_R/2)(333/106)-w\).  Positivity therefore proves both strict
half-period location rows without a decimal approximation.
\begin{center}
\small
\begin{tabular}{c|c|c|c|c}
\toprule
\((\ell,n)\)&\(w\)&\((h_L,h_R)\)&left margin&right margin\\
\midrule
\((1,2)\)&\(77/10\)&\((4,5)\)&\(99/70\)&\(163/1060\)\\
\((2,2)\)&\(9\)&\((5,6)\)&\(8/7\)&\(45/106\)\\
\((3,2)\)&\(103/10\)&\((6,7)\)&\(61/70\)&\(737/1060\)\\
\((5,2)\)&\(129/10\)&\((8,9)\)&\(23/70\)&\(1311/1060\)\\
\((1,3)\)&\(21/2\)&\((6,7)\)&\(15/14\)&\(105/212\)\\
\((2,3)\)&\(61/5\)&\((7,8)\)&\(6/5\)&\(97/265\)\\
\((6,1)\)&\(10\)&\((6,7)\)&\(4/7\)&\(211/212\)\\
\((7,1)\)&\(23/2\)&\((7,8)\)&\(1/2\)&\(113/106\)\\
\bottomrule
\end{tabular}
\end{center}

\subsection{Sixteen half-period signs}

At an integral or half-integral multiple of \(\pi\), exactly one of
\(\sin x,\cos x\) vanishes.  Direct substitution in the displayed
polynomials gives the following complete sign table.
\begin{center}
\small
\begin{tabular}{c|c|c@{\qquad}c|c|c}
\toprule
\((\ell,n)\)&\(h_L\)&left sign&\((\ell,n)\)&\(h_R\)&right sign\\
\midrule
\((1,2)\)&4&\(-\)&\((1,2)\)&5&\(+\)\\
\((2,2)\)&5&\(-\)&\((2,2)\)&6&\(+\)\\
\((3,2)\)&6&\(-\)&\((3,2)\)&7&\(+\)\\
\((5,2)\)&8&\(-\)&\((5,2)\)&9&\(+\)\\
\((1,3)\)&6&\(+\)&\((1,3)\)&7&\(-\)\\
\((2,3)\)&7&\(+\)&\((2,3)\)&8&\(-\)\\
\((6,1)\)&6&\(+\)&\((6,1)\)&7&\(-\)\\
\((7,1)\)&7&\(+\)&\((7,1)\)&8&\(-\)\\
\bottomrule
\end{tabular}
\end{center}
Here are explicit polynomial witnesses for the only signs not immediate
from \(A_1=1,C_1=-1,A_2=3-t,C_2=-3,A_3=15-6t,C_3=t-15\).
From \eqref{eq:pi-bounds}, at \(x=4\pi\) one has
\(144<t<159\), at \(x=9\pi/2\) one has \(t>182\), and at
\(x=7\pi/2\) one has \(110<t<121\).  Monotonicity on these rational
intervals and direct evaluation give
\[
 C_5(t)<C_5(144)=-6561,\qquad
 A_5(t)>A_5(182)=421365,
\]
\[
 C_6(t)<C_6(81)=-46116,\qquad
 A_6(t)>A_6(110)=700645,
\]
\[
 A_7(t)<A_7(110)=-5878565,\qquad
 C_7(t)<C_7(144)=-2492559.
\]
For example, \(A_6'=-4725+420t-3t^2\) is positive on
\([110,121]\), while \(A_7'\) is negative there; the other asserted
monotonicities follow at once from their displayed derivatives.  Combining
these signs with the quadrant signs of sine and cosine proves all sixteen
entries.

\subsection{Eight rational-wall signs}

Substitution in \eqref{eq:P-rec} reduces each rational-wall sign to one
of the following exact tangent comparisons.  The shifted angle is always
in \((0,\pi/2)\), as already proved by the location table.
\begin{center}
\small
\begin{tabular}{c|c|p{0.64\linewidth}|c}
\toprule
\((\ell,n)\)&wall&hand certificate&\(\sgn\mathsf j_\ell(w)\)\\
\midrule
\((1,2)\)&\(77/10\)&For \(u=5\pi/2-w\),
 \(\tan u>3/20>10/77=1/w\); hence
 \(\tan(w-2\pi)=\cot u<w\).&\(-\)\\
\((2,2)\)&\(9\)&For \(u=3\pi-w\),
 \(\tan u>21/50>9/26=3w/(w^2-3)\).&\(-\)\\
\((3,2)\)&\(103/10\)&For \(u=7\pi/2-w\),
 \(\tan u>69/100>621540/938227=(6w^2-15)/(w(w^2-15))\).&\(-\)\\
\((5,2)\)&\(129/10\)&For \(y=w-4\pi\), \(0<y<17/50\), so
 \(\tan y<1700/4711<177800829/427700150=-wC_5/A_5\).&\(-\)\\
\((1,3)\)&\(21/2\)&For \(u=7\pi/2-w\),
 \(\tan u>49/100>2/21=1/w\), hence
 \(\tan(w-3\pi)<w\).&\(+\)\\
\((2,3)\)&\(61/5\)&For \(u=4\pi-w\),
 \(\tan u>9/25>915/3646=3w/(w^2-3)\).&\(+\)\\
\((6,1)\)&\(10\)&For \(y=w-3\pi\),
 \(4/7<y<29/50\), so
 \(\tan y<2900/4159<188790/127579=-wC_6/A_6\).&\(+\)\\
\((7,1)\)&\(23/2\)&For \(u=4\pi-w\),
 \(\tan u>u+u^3/3>546377/375000>
 3153151489/2276518664=wC_7/A_7\).&\(+\)\\
\bottomrule
\end{tabular}
\end{center}
For complete cross-multiplication visibility, the four less immediate
positive gaps in that table are
\begin{align*}
 \frac{69}{100}-\frac{621540}{938227}
 &=\frac{2583663}{93822700},\\
 \frac{177800829}{427700150}-\frac{1700}{4711}
 &=\frac{110529450419}{2014895406650},\\
 \frac{188790}{127579}-\frac{2900}{4159}
 &=\frac{415198510}{530601061},\\
 \frac{546377}{375000}-\frac{3153151489}{2276518664}
 &=\frac{3837851856583}{53355906187500}.
\end{align*}
The four small gaps are respectively
\(31/1540,24/325,829/2100,9939/91150\).

\subsection{Intermediate rational-wall signs used by the index induction}

The coefficient signs make all unmentioned entries in the sign vectors
immediate.  The remaining intermediate entries require the following exact
thresholds.  At \(w=10\), with \(y=10-3\pi\),
\[
 \tan y<\frac{2900}{4159}<\frac{170}{117}<10,
 \qquad
 \tan y>y>\frac47>\frac{890}{21789}.
\]
Substitution in \(P_1,P_3,P_5\) gives
\(\sgn(\mathsf j_1,\mathsf j_3,\mathsf j_5)=(+,-,-)\), and
\(P_4=7P_3-w^2P_2<0\).  At \(w=103/10\), with
\(u=7\pi/2-w\),
\[
 \tan u>\frac{69}{100}>\frac{10}{103},
\]
which gives \(\mathsf j_1(w)>0\).  At \(w=23/2\), with \(u=4\pi-w\),
\[
 \tan u>\frac{546377}{375000}>
 \frac{3153151489}{2276518664}>
 \frac{224020}{186301}>\frac{138}{517}.
\]
Besides the target \(\mathsf j_7\) sign, the last two thresholds give
\(\mathsf j_4(w)<0\) and \(\mathsf j_2(w)>0\); the recurrence gives the
intervening signs.  Finally, at \(w=129/10\), with \(y=w-4\pi\),
\[
 \tan y<\frac{1700}{4711}<
 \frac{177800829}{427700150}<
 \frac{651063}{327820}<\frac{129}{10}.
\]
Besides the target \(\mathsf j_5\) sign, this proves
\(\mathsf j_3(w)>0\) and \(\mathsf j_1(w)<0\), and the recurrence gives
the remaining signs.  Every sign used below is therefore supported by a
displayed rational comparison.

\subsection{Exact zero indices from adjacent-order interlacing}

For \(x>0\), define the strict positive-zero count
\[
 Z_\ell(x):=\#\{n\in\mathbb N:j_{\ell+1/2,n}<x\}.
\]

\begin{lemma}[Interlacing-to-index principle]
\label{lem:zero-index}
For every \(\ell\in\mathbb N_0\) and \(n\in\mathbb N\),
\begin{equation}
 j_{\ell+1/2,n}<j_{\ell+3/2,n}<j_{\ell+1/2,n+1}.
 \label{eq:adjacent-interlacing}
\end{equation}
If \(x\) is not a zero of
\(\mathsf j_0,\ldots,\mathsf j_L\), then, for \(0\leq\ell<L\),
\begin{equation}
 Z_{\ell+1}(x)\in
 \bigl\{\max\{Z_\ell(x)-1,0\},\,Z_\ell(x)\bigr\},
 \qquad
 \sgn\mathsf j_{\ell+1}(x)=(-1)^{Z_{\ell+1}(x)}.
 \label{eq:index-recursion}
\end{equation}
Consequently \(Z_0(x)\), together with the signs of
\(\mathsf j_1(x),\ldots,\mathsf j_L(x)\), determines all the counts
\(Z_1(x),\ldots,Z_L(x)\) uniquely.
\end{lemma}

\begin{proof}
DLMF~10.21.2 states, for \(\nu\geq0\),
\[
 j_{\nu,1}<j_{\nu+1,1}<j_{\nu,2}<j_{\nu+1,2}
 <j_{\nu,3}<\cdots .
\]
Taking \(\nu=\ell+1/2\) proves
\eqref{eq:adjacent-interlacing}.  Put \(q=Z_\ell(x)\).  The first
inequality in \eqref{eq:adjacent-interlacing} shows that at most \(q\)
zeros of order \(\ell+3/2\) lie below \(x\).  If \(q\geq1\), the second
inequality shows that its first \(q-1\) zeros lie below
\(j_{\ell+1/2,q}<x\).  Thus
\[
 \max\{q-1,0\}\leq Z_{\ell+1}(x)\leq q,
\]
which is the first assertion in \eqref{eq:index-recursion}.

The spherical Bessel function \(\mathsf j_\ell\) is positive for all
sufficiently small positive arguments.  Its positive zeros are simple by
uniqueness for the Bessel equation, so its sign changes at every positive
zero.  This proves the parity identity in
\eqref{eq:index-recursion}.  When \(q\geq1\), the two possible counts have
opposite parity; when \(q=0\), only zero is possible.  The sign therefore
selects the count uniquely.
\end{proof}

Because \(\mathsf j_0(x)=\sin x/x\), one has
\(j_{1/2,n}=n\pi\).  The strict locations in the first table therefore
give \(Z_0(w)\) exactly.  Applying Lemma~\ref{lem:zero-index} successively
to the nonzero sign vectors obtained from \eqref{eq:P-rec} and the exact
tangent comparisons above gives:
\begin{center}
\scriptsize
\renewcommand{\arraystretch}{1.16}
\begin{tabular}{c|c|c|c|c|c}
\toprule
wall \(w\)&target \((\ell,n)\)&location of \(w\)&\(Z_0(w)\)&
\(\sgn(\mathsf j_0,\ldots,\mathsf j_\ell)(w)\)&
\((Z_0,\ldots,Z_\ell)(w)\)\\
\midrule
\(77/10\)&\((1,2)\)&\((2\pi,5\pi/2)\)&2&
\((+,-)\)&\((2,1)\)\\
\(9\)&\((2,2)\)&\((5\pi/2,3\pi)\)&2&
\((+,+,-)\)&\((2,2,1)\)\\
\(10\)&\((6,1)\)&\((3\pi,7\pi/2)\)&3&
\((-,+,+,-,-,-,+)\)&\((3,2,2,1,1,1,0)\)\\
\(103/10\)&\((3,2)\)&\((3\pi,7\pi/2)\)&3&
\((-,+,+,-)\)&\((3,2,2,1)\)\\
\(21/2\)&\((1,3)\)&\((3\pi,7\pi/2)\)&3&
\((-,+)\)&\((3,2)\)\\
\(23/2\)&\((7,1)\)&\((7\pi/2,4\pi)\)&3&
\((-,-,+,+,-,-,-,+)\)&\((3,3,2,2,1,1,1,0)\)\\
\(61/5\)&\((2,3)\)&\((7\pi/2,4\pi)\)&3&
\((-,-,+)\)&\((3,3,2)\)\\
\(129/10\)&\((5,2)\)&\((4\pi,9\pi/2)\)&4&
\((+,-,-,+,+,-)\)&\((4,3,3,2,2,1)\)\\
\bottomrule
\end{tabular}
\end{center}

In every row the target component is \(Z_\ell(w)=n-1\), and the displayed
sign is nonzero.  Hence, with the convention \(j_{\nu,0}=0\),
\[
 j_{\ell+1/2,n-1}<w<j_{\ell+1/2,n}.
\]
Thus the unique next zero is the claimed \(n\)-th zero, and the eight rows
prove
\begin{equation}
\begin{gathered}
 j_{3/2,2}>77/10,\qquad j_{5/2,2}>9,\qquad
 j_{7/2,2}>103/10,\qquad j_{11/2,2}>129/10,\\
 j_{3/2,3}>21/2,\qquad j_{5/2,3}>61/5,\qquad
 j_{13/2,1}>10,\qquad j_{15/2,1}>23/2.
\end{gathered}
\label{eq:indexed-zero-registry}
\end{equation}

\subsection{Six Lorch rows and two wall-order rows}

For \(\nu=\ell+1/2\), the second strict inequality, item~(ii), displayed
in the published abstract of Lorch's article \emph{Some inequalities for
the first positive zeros of Bessel functions}, SIAM J.\ Math.\ Anal.\
\textbf{24} (1993), 814--823, doi:10.1137/0524050, is valid for
\(-1<\nu<\infty\) and states
\[
 j_{\nu,1}^2>
 \frac{24(\nu+1)^2}{1-2\nu+\sqrt{(2\nu+3)(2\nu+11)}}
 -2(\nu^2-1)=:R_\ell.
\]
After rationalization,
\[
 R_\ell=\frac{3(2\ell+3)\sqrt{(\ell+2)(\ell+6)}
 -2\ell^2+\ell+6}{4}.
\]
All sides being squared are positive, and the following integer residuals
prove the three required strict specializations:
\begin{center}
\begin{tabular}{c|c|c}
\toprule
\(\ell\)&wall \(c\)&positive squared residual proving \(R_\ell>c^2\)\\
\midrule
8&\(71/6\)&\(1026^2\cdot35-6067^2=35171\)\\
9&\(64/5\)&\(1575^2\cdot165-20059^2=6939644\)\\
10&\(69/5\)&\(3450^2\cdot3-5911^2=767579\)\\
\bottomrule
\end{tabular}
\end{center}
Since \(\mathsf j_\ell(x)>0\) for sufficiently small \(x>0\), and its
first positive zero is strictly to the right of \(c\), these same three
rows imply \(\mathsf j_\ell(c)>0\).  Thus the three executable sign rows
are analytic corollaries, not independent numerical assumptions.  Finally,
\[
 \frac{103}{10}-\frac{81}{8}=\frac7{40}>0,
 \qquad
 \frac{61}{5}-\frac{21}{2}=\frac{17}{10}>0.
\]

The one-to-one zero registry is as follows.  Number the eight internal
faces in the order of the location table.  For face \(r\), IDs
\(Z_{5r-4},\ldots,Z_{5r}\) are, in order: right of the left half-period,
left of the right half-period, left half-period sign, right half-period
sign, and rational-wall sign.  This bijects IDs \(Z1\)--\(Z40\) with the
first forty source rows.  The remaining bijection is
\begin{center}
\small
\begin{tabular}{c|c|l}
\toprule
ID&source row&predicate\\
\midrule
Z41&41&Lorch specialization \(\ell=8,c=71/6\)\\
Z42&42&positive sign at \(\ell=8,c=71/6\)\\
Z43&43&Lorch specialization \(\ell=9,c=64/5\)\\
Z44&44&positive sign at \(\ell=9,c=64/5\)\\
Z45&45&Lorch specialization \(\ell=10,c=69/5\)\\
Z46&46&positive sign at \(\ell=10,c=69/5\)\\
Z47&47&\(103/10>81/8\)\\
Z48&48&\(61/5>21/2\)\\
\bottomrule
\end{tabular}
\end{center}

\section{The forty-three lower-closure predicates}

\subsection{Eight activation and exclusion rows}

The complete arithmetic is
\begin{align*}
 (23/2)^2-114&=73/4>0,&10^2&=100,\\
 (129/10)^2-114&=5241/100>0,
 &103/10-81/8&=7/40>0,\\
 (81/8)^2+8&=7073/64,
 &(81/8)^2+18-114&=417/64>0,\\
 (21/2)^2+4-114&=1/4>0,
 &16\pi^2&>16\cdot3^2>114.
\end{align*}
These are exactly lower IDs L1--L8, corresponding to source rows 1--8.

\subsection{Five Weyl payments}

For \(0<\rho<\rc<1/7\), \eqref{eq:pi-bounds} yields the uniform analytic
lower bound
\begin{equation}
 W(\rho,K)=\frac{2}{9\pi}(1-\rho^3)K^3
 >\frac{38}{539}K^3.
 \label{eq:weyl-coarse}
\end{equation}
Using \(d>9\) and \(\sqrt{7073}/8>21/2\), the five payments are the
following exact positive reserves:
\begin{center}
\begin{tabular}{c|c|c|c}
\toprule
ID&source row&threshold/cap&reserve from \eqref{eq:weyl-coarse}\\
\midrule
L9&10&\(K>d,\ 46\)&\((38/539)9^3-46=2908/539\)\\
L10&11&\(K>10,\ 59\)&\((38/539)10^3-59=6199/539\)\\
L11&12&\(K>81/8,\ 66\)&\((38/539)(81/8)^3-66=990435/137984\)\\
L12&13&\(K>21/2,\ 69\)&\((38/539)(21/2)^3-69=555/44\)\\
L13&14&\(K>\sqrt{7073}/8,\ 78\)&\((38/539)(21/2)^3-78=159/44\)\\
\bottomrule
\end{tabular}
\end{center}

\subsection{The cap--395 cell: eighteen rows}

Put \(R=367/20\), \(z_\ell=\ell+1/2\).  The analytic turning estimate
\[
 G_K(z)<\frac{(K-z)^{3/2}}{3\sqrt K}
\]
is increasing in \(K\), so on \(K<R\) a strict cap \(c_\ell\) follows
from
\begin{equation}
 \frac{(R-z_\ell)^3}{9R}
 =\frac{(357-20\ell)^3}{1321200}
 <\left(c_\ell+\frac34\right)^2.
 \label{eq:proxy-square}
\end{equation}
The following table gives all fifteen caps and the exact positive numerator
\(82575(4c_\ell+3)^2-(357-20\ell)^3\) obtained by clearing the
positive denominators in \eqref{eq:proxy-square}.
\begin{center}
\scriptsize
\begin{tabular}{c|rrrrrrrr}
\toprule
\(\ell\)&0&1&2&3&4&5&6&7\\
\midrule
\(c_\ell\)&6&5&5&4&4&3&3&3\\
margin
&14697882&5409422&11827162&3611502&8555642
&1604782&5267322&8361062\\
\bottomrule
\end{tabular}

\medskip
\begin{tabular}{c|rrrrrrr}
\toprule
\(\ell\)&8&9&10&11&12&13&14\\
\midrule
\(c_\ell\)&2&2&1&1&1&1&0\\
margin&2346202&4446342&176282&1474822&2444562&3133502&286642\\
\bottomrule
\end{tabular}
\end{center}
Thus L15--L29 correspond bijectively to source rows 16--30, with
L\((15+\ell)\) the \(\ell\)-th table column.  Since \(G_R(z)\) is
strictly decreasing in \(z\), the \(c_{14}=0\) row removes every
\(\ell\ge14\) tail; this is L30/source row 31.

The weighted cap itself, L14/source row 15, is visible from the contributions
\[
 6,15,25,28,36,33,39,45,34,38,21,23,25,27,
 \qquad \sum_{\ell=0}^{13}(2\ell+1)c_\ell=395.
\]
Finally, on the exceptional cell \(\rho K\ge5/2\), with
\(\rho<11/80\) and \(1/\pi>113/355\),
\begin{align*}
 W(\rho,K)
 &>\frac{2}{9}\frac{113}{355}\left(\frac52\right)^3
 \left[\left(\frac{80}{11}\right)^3-1\right]\\
 &=\frac{160293325}{378004}
 =395+\frac{10981745}{378004}>395.
\end{align*}
This is L31/source row 32.

\subsection{Twelve floor and boundary-owner rows}

The nine finite arithmetic rows are instances of one identity:
\begin{equation}
 \frac32(p+1)\le 2p+\frac12
 \quad(p\ge2),\qquad
 \left(2p+\frac12\right)-\frac32(p+1)=\frac{p-2}{2}\ge0.
 \label{eq:floor-general}
\end{equation}
At \(p=2\), \((3/2)3=9/2\); this is L32/source row 33.
For \(p=2,\ldots,9\), formula \eqref{eq:floor-general} is respectively
L33--L40/source rows 34--41.

The half-integer face is exact:
\(5/2=2+1/2\).  Assigning it to the high side makes the interface
remainder vanish and leaves the exceptional cell closed; this is
L41/source row 42.  At \(K=2/\om\),
\[
 p=\lfloor\om K\rfloor=2,
 \qquad \rho K\le\frac{2\rc}{\om}<3
\]
by T10.  This proves L42/source row 43.  Moreover
\[
 \rho K-\frac12<\frac52,\qquad
 m=\left\lfloor\rho K-\frac12\right\rfloor\le2
 \le 2p-1=3,
\]
and the strict exceptional condition \(K<2/\om\) fails at equality.
Thus the complementary average-floor branch owns the equality face,
proving L43/source row 44.
More generally, if \(K\ge2/\om\), then \(p\ge2\) and
\[
 \rho K<\frac32(p+1)\le2p+\frac12,
 \qquad m\le2p-1,
\]
so the finite instances conceal no unproved large-\(p\) case.

\section{The supplemental cap--74 incompatibility}

The propagated ball wall requires \(z^2>2409/100\), while the product
wall requires \(z^2<70/3\).  They are incompatible because
\[
 \frac{2409}{100}-\frac{70}{3}
 =\frac{7227-7000}{300}=\frac{227}{300}>0.
\]
This is supplemental ID S1 and is a hand proof, not merely a reported
certificate value.

\section{Census and source-visibility verdict}

\begin{center}
\begin{tabular}{l|r|l|l}
\toprule
family&count&IDs&hand-visible status before this note\\
\midrule
threshold remainder&11&T1--T11&printed, with a few implications compressed\\
zero specializations&48&Z1--Z48&tables printed; several recurrence steps compressed\\
lower remainder&43&L1--L43&payments printed; floor instances generated implicitly\\
supplemental cap--74&1&S1&printed exactly\\
\midrule
total&103&&\\
\bottomrule
\end{tabular}
\end{center}

No predicate in this 103-row allocation is an unsupported executable-only
mathematical premise.  Before this note, the principal visibility gaps were
the one-to-one expansion of the forty zero sign/location rows, the fact that
the three Lorch-wall signs are logical corollaries of the strict first-zero
bounds, and the expansion of the twelve floor/boundary rows.  The displayed
recurrences, rational margins, general floor identity, and exact ownership
arguments close those presentation gaps.

\end{document}
